\section{Managerial and policy implications}

\subsection{Recommendations should expose their decision layer}

A crop-ranking tool can remain useful without prescribing acreage. The problem
arises when an ordinal output is presented as if the top crop should receive
the most land. A prescriptive system should expose at least five objects: the
cardinal margin model, the joint distribution used to form portfolio losses,
the operational feasible set, the risk convention and the selector used when
optima are multiple. If these objects are unavailable, the defensible output
is a rank or scenario comparison, not a claim of optimal acreage.

The distinction matters for accountability. A lender or platform that flags a
farmer for planting less of a high-scoring crop may be penalizing compliance
with a rotation, contract or liquidity restriction omitted from the score.
Conversely, saying that every disagreement is rational would make the
recommendation system unfalsifiable. The remedy is to compare declared
feasible counterfactuals and audit the optimal face.

\subsection{Downside-risk governance}

The Kansas experiment shows the cost of treating a mean--variance solution as
tail safe. Under matched rank dependence, the conventional policy earns a
higher mean but violates the true-law CVaR ceiling. This suggests a governance
test for decision-support systems: report the allocation under at least one
heavy-tailed joint stress law, evaluate it under the same downside-risk
metric used for the farmer's tolerance, and disclose whether the policy
remains feasible. The test is more informative than reporting a diversification
index or pairwise correlation alone.

The result does not justify universal conservatism. A very tight risk ceiling
reduces expected profit and can alter which crop dominates; a loose ceiling
may leave the expected-profit allocation unchanged. Decision makers should
therefore report the entire risk frontier and active-set transitions. The
location of reversal is conditional on that frontier, not an immutable
property of a crop.

\subsection{Invest in information only with an action pathway}

Forecast accuracy and operational flexibility should be evaluated jointly.
When a share of acreage can be adjusted after a credible pre-plant signal,
better information can be increasingly valuable. If planting is already
committed or posterior-optimal plans coincide, more signal accuracy has zero
economic value even if the prediction is statistically informative. These
are design criteria for climate services: document decision timing and the
specific acreage, input, irrigation or harvest actions that remain open.

Shock-buffering flexibility creates a different investment trade-off.
Irrigation, robust varieties, input switching or flexible harvest capacity can
reduce exposure to the predicted state. This may raise total profit while
reducing the incremental value of a forecast. Substitution is not a failure of
the forecast; it means the operational technology has made the distinction
between states less consequential. A portfolio of information and flexibility
should therefore be assessed by total optimized value, not by requiring the
interaction to be positive.

\subsection{Policy interpretation}

Rotation and diversification mandates should distinguish crop count from
joint-tail protection. A policy that increases the number of planted crops can
still leave a producer exposed to simultaneous drought losses when the crops
share phenological vulnerabilities \citep{schmitt2024drought}. Crop-pair
choice, timing, geography and access to off-farm risk instruments may matter
more than a generic diversity threshold.

Likewise, public rankings based on climate suitability should not be evaluated
solely against observed acreage. State acreage aggregates many farms and
combines agronomic, financial and institutional decisions. The appropriate
policy question is whether the ranking improves a declared decision under
realistic constraints. The external panel in this study motivates that
question but cannot answer it causally.
